\begin{helper}[Comparison principle on a ball]
Let $n\ge2$, $p>1$, $\varepsilon\ge0$ and $0<\gamma_0\le\gamma_1$, and suppose
the Sobolev--Poincar\'e inequality on balls holds in dimension $n$
(\noderef{CitedSobolevPoincareBall}).  Let
$A\colon\R^n\to\R^n$ satisfy (Q1) of the thesis--(Q2) of the thesis with these parameters
on all of $\R^n$, i.e.\ \noderef{EllipticStructureQ12} holds with
$\Omega=\R^n$ for the field $(x,z,h)\mapsto A(h)$, which is the setting of \S3
where $A$ depends on $h$ alone.  Let $B_R=B_R(x_0)$ with $R>0$ and let $u,v\in
W^{1,p}(B_R)$ (\noderef{MemW1pOn}) be weak solutions of
$\operatorname{div}A(\nabla u)=0$ in $B_R$ (\noderef{IsWeakSolutionOn}), with
weak gradients $g_u$, $g_v$.  Suppose $U$ and $V$ are continuous on
$\overline B_R$ and extend $u$ and $v$ from $B_R$, and suppose
\[
U\le V\qquad\text{on }\partial B_R .
\]
Then $u\le v$ almost everywhere in $B_R$.

The hypothesis \noderef{CitedSobolevPoincareBall} is the cited result (C6) of the
paper's Note for the formalization, carried here as an explicit argument because
the proof uses it: it is what turns ``$\varphi_\eta\in W^{1,1}(B_R)$ with weak
gradient $0$'' into ``$\varphi_\eta$ equals its own mean on $B_R$ almost
everywhere'' in Step 4.  It is a statement about $n$ alone, and $n$ is bound
first, so it introduces no dependence on any quantity outside the list this
lemma already names and carries no numerical content of its own.
\end{helper}
\begin{proof}
Write $B_r=B_r(x_0)$ and $S_r=\partial B_r(x_0)$ throughout, and set
$w=u-v$ on $B_R$, $W=U-V$ on $\overline B_R$.  Note $0<|B_R|<\infty$, since
$B_R$ is a ball of positive finite radius.

\emph{Step 1: strict monotonicity of $A$.}  Put $r_\ast=\gamma_1/\gamma_0$,
which is well defined and positive because $\gamma_0>0$ and
$\gamma_1\ge\gamma_0>0$, and note $\gamma_1=r_\ast\gamma_0$, so in particular
$\gamma_1\le r_\ast\gamma_0$.  Let $C=C(n,p)>0$ be the constant supplied by
\noderef{StructureMonotonicityBound} for this $n$, $p$ and $r_\ast$, available
because $n\ge2$, $p>1$ and $r_\ast>0$.  Apply that node with the present
$\varepsilon,\gamma_0,\gamma_1$, with $\Omega=\R^n$ and with the field
$(x,z,h)\mapsto A(h)$: its remaining hypotheses are $\varepsilon\ge0$,
$\gamma_0>0$, $\gamma_1\le r_\ast\gamma_0$ and
\noderef{EllipticStructureQ12} for that field on $\R^n$, all of which are
hypotheses of the present lemma or were just verified.  Its conclusion is that
for almost every $x\in\R^n$, and then for every $z\in\R$ and all
$a,b\in\R^n$,
\[
C\,\gamma_0\bigl(\varepsilon+|a|^2+|b|^2\bigr)^{\frac p2-1}|a-b|^2
\;\le\;\bigl(A(a)-A(b),\,a-b\bigr)
\]
(the membership $x\in\Omega=\R^n$ being automatic).  The displayed inequality
involves neither $x$ nor $z$, because the field $(x,z,h)\mapsto A(h)$ does not;
and the set of admissible $x$ has full Lebesgue measure in $\R^n$, hence is
nonempty, Lebesgue measure on $\R^n$ not being the zero measure.  Fixing one
admissible $x$ and any $z$ therefore yields, for all $a,b\in\R^n$,
\[
\bigl(A(a)-A(b),\,a-b\bigr)
\;\ge\;C\,\gamma_0\,\bigl(\varepsilon+|a|^2+|b|^2\bigr)^{\frac p2-1}\,|a-b|^2 .
\tag{1}
\]

Note what (1) does and does not say.  The right-hand side is nonnegative
always, so $A$ is monotone.  We will \emph{not} assert that the weight is
everywhere positive; what we need is only the implication
\[
\bigl(A(a)-A(b),a-b\bigr)=0\ \Longrightarrow\ a=b .
\tag{2}
\]
To get (2), suppose $(A(a)-A(b),a-b)=0$.  Then (1) gives
$C\gamma_0(\varepsilon+|a|^2+|b|^2)^{\frac p2-1}|a-b|^2\le0$, and $C\gamma_0>0$,
so
\[
\bigl(\varepsilon+|a|^2+|b|^2\bigr)^{\frac p2-1}\,|a-b|^2\le0 .
\]
Both factors are nonnegative --- the first is a real power of the nonnegative
real $\varepsilon+|a|^2+|b|^2$, the second a square --- so their product is $0$,
and hence one of them is $0$.  If $|a-b|^2=0$ we are done.  Otherwise
$(\varepsilon+|a|^2+|b|^2)^{\frac p2-1}=0$.  A real power $t^{s}$ of a
nonnegative real $t$ vanishes only when $t=0$ and $s\ne0$: for $t>0$ one has
$t^{s}>0$, and for $t=0$ one has $0^{0}=1$.  So this second alternative forces
$\varepsilon+|a|^2+|b|^2=0$, i.e.\ $\varepsilon=0$ and $a=b=0$, in which case
$a=b$ as well.  This proves (2).

\emph{Step 2: the test functions.}  Fix $\eta>0$ and let $\varphi_\eta$ be the
function on $\R^n$ given by
\[
\varphi_\eta=\begin{cases}
(w-\eta)^+=\max\{u-v-\eta,\,0\} & \text{on }B_R,\\
0 & \text{off }B_R .
\end{cases}
\]
(The cut-off off $B_R$ matters only for the support bookkeeping below; on $B_R$,
where the two weak formulations live, $\varphi_\eta$ is just $(w-\eta)^+$.)
We check that $\varphi_\eta$ is an admissible test function for both weak
formulations, i.e.\ that it lies in $W^{1,p}(B_R)$ with compact support contained
in $B_R$, which is exactly the test class of \noderef{IsWeakSolutionOn}.

\emph{Support.}  Put $K_\eta=\{x\in\overline B_R:W(x)\ge\eta\}$.  Since $U$ and
$V$ are continuous on $\overline B_R$, so is $W$, hence $K_\eta$ is closed in the
compact set $\overline B_R$ and therefore compact.  Moreover
$K_\eta\cap S_R=\emptyset$: on $S_R$ we have $U\le V$ by hypothesis, i.e.\
$W\le0<\eta$.  Hence $K_\eta\subseteq\overline B_R\setminus S_R=B_R$, the set
identity being that a point of the closed ball is at distance $\le R$ from
$x_0$, and is in the sphere exactly when that distance equals $R$, so removing
the sphere leaves exactly the points at distance $<R$, which is $B_R$.  Now if
$\varphi_\eta(x)\ne0$ for
some $x\in B_R$ then $w(x)>\eta$, and $W(x)=U(x)-V(x)=u(x)-v(x)=w(x)>\eta$
because $U=u$ and $V=v$ on $B_R$; so $x\in K_\eta$.  Since $\varphi_\eta$
vanishes off $B_R$ by construction, this gives
$\{\varphi_\eta\ne0\}\subseteq K_\eta$ on all of $\R^n$, and since $K_\eta$ is
closed the closure of that set --- the support of $\varphi_\eta$ --- is contained
in $K_\eta$.  So $\varphi_\eta$ has compact support and
$\operatorname{supp}\varphi_\eta\subseteq B_R$, which is exactly the test class
of \noderef{IsWeakSolutionOn}.

\emph{Sobolev regularity.}  By hypothesis $u,v\in W^{1,p}(B_R)$ with weak
gradients $g_u,g_v$ (\noderef{MemW1pOn}).  We first record that
$w=u-v\in W^{1,p}(B_R)$ with weak gradient $g_u-g_v$, writing the splitting of
the integrals out rather than appealing to linearity in the abstract.  Let
$\varphi$ be smooth with compact support contained in $B_R$ and let
$i\in\{1,\dots,n\}$.  By the $L^p$ clause of \noderef{MemW1pOn},
$u,v\in L^p(B_R)$ and $|g_u|,|g_v|\in L^p(B_R)$; since $|B_R|<\infty$ and
$p\ge1$, H\"older on a set of finite measure (Preamble) puts all four of
$u,v,g_u^i,g_v^i$ in $L^1(B_R)$.  The functions $\varphi$ and $D_i\varphi$ are
continuous with compact support, hence bounded on $\R^n$, and the product of an
$L^1(B_R)$ function $f$ with a function bounded by $\Lambda$ is in $L^1(B_R)$,
being measurable and dominated by $\Lambda|f|\in L^1(B_R)$ (domination,
Preamble).  So
each of the four integrals
\[
\int_{B_R}u\,D_i\varphi,\quad
\int_{B_R}v\,D_i\varphi,\quad
\int_{B_R}g_u^i\,\varphi,\quad
\int_{B_R}g_v^i\,\varphi
\]
is finite, and subtracting the defining identity of
\noderef{HasWeakGradientOn} for $(v,g_v)$ from the one for $(u,g_u)$ is
legitimate and gives
\[
\int_{B_R}(u-v)\,D_i\varphi=-\int_{B_R}\bigl(g_u^i-g_v^i\bigr)\varphi ,
\]
which is the defining identity of \noderef{HasWeakGradientOn} for the pair
$(w,g_u-g_v)$ on $B_R$.  Moreover $w\in L^p(B_R)$ and $|g_u-g_v|\in L^p(B_R)$
by the $L^p$ triangle inequality (Minkowski, Preamble).  Hence
$w\in W^{1,p}(B_R)$ with weak gradient $g_u-g_v$ (\noderef{MemW1pOn}).

Now $B_R$ is open with $|B_R|<\infty$ and $p\ge1$, so
\noderef{WeakGradientPositivePart} applied to
$w$ with $c=\eta$ gives
\[
\varphi_\eta=(w-\eta)^+\in W^{1,p}(B_R)\quad\text{with weak gradient}\quad
g_\eta=(g_u-g_v)\,\chi_{\{w>\eta\}} ,
\tag{3}
\]
the $L^p$ memberships being part of that conclusion.  (Off $B_R$ our $\varphi_\eta$ is cut off to $0$.  This affects neither half of
the conclusion (3): not the weak gradient, since
\noderef{HasWeakGradientOn} on $B_R$ only ever integrates against test functions
supported in $B_R$ and so sees no value of $\varphi_\eta$ outside $B_R$; and not
the two $L^p$ memberships, since those are memberships with respect to the
measure restricted to $B_R$, which likewise sees no value outside $B_R$.)

\emph{Step 3: testing, integrability, and subtracting.}

\emph{(3a) Measurability.}  By the $L^p$ clause of \noderef{MemW1pOn} --- its
item 2, which carries both memberships --- the functions $u,v$ and the fields
$g_u,g_v$ are measurable on $B_R$, hence so are $w=u-v$ and $g_u-g_v$.  (That
item states plain measurability, which is what is used here and below; no
almost-everywhere qualification is needed or claimed.)
By the \emph{continuity} clause of \noderef{EllipticStructureQ12}
(clause (1), read with $\Omega=\R^n$) the map $h\mapsto A(h)$ is continuous on
all of $\R^n$, so $A\circ g_u$ and $A\circ g_v$ are measurable on $B_R$.  It is
clause (1) and not the differentiability clause that is cited here, and the
distinction is not cosmetic: clause (2) gives $C^1$ only on
$D(\varepsilon)=\{h:\varepsilon+|h|^2\ne0\}$, and the present statement allows
$\varepsilon\ge0$, so at $\varepsilon=0$ it says nothing at $h=0$.  An
admissible field need not be differentiable there, nor even Lipschitz: for
$1<p<2$ the field $A(h)=|h|^{p-2}h$ satisfies (Q1) of the thesis--(Q2) of the thesis at
$\varepsilon=0$ with $\gamma_0=\gamma_1=p-1>0$, while
$|A(h)-A(0)|/|h|=|h|^{p-2}\to\infty$ as $h\to0$.  Continuity on all of $\R^n$
is what clause (1) asserts and is exactly and only what measurability of
$A\circ g_u$ and $A\circ g_v$ needs.  The
superlevel set $\{w>\eta\}$ of the measurable function $w$ is measurable
(Preamble), so its intersection with the open set $B_R$ is measurable, and
therefore $g_\eta=(g_u-g_v)\chi_{\{w>\eta\}}$ is measurable and so is the
scalar function
\[
x\longmapsto\bigl(A(g_u(x))-A(g_v(x)),\,g_u(x)-g_v(x)\bigr),
\]
the Euclidean inner product being continuous.  These two measurability facts are
what the ``nonnegative integral zero'' inference at the end of this step
presupposes.

\emph{(3b) Integrability, via the growth bound.}  This is where (Q2) of the thesis
enters: (Q1) of the thesis alone says nothing about the size of $A$.  Let
$c_\ast=c_\ast(n,p)>0$ be the constant of \noderef{StructureGrowthBound} for
this $n$ and $p$, and apply that node with the present
$\varepsilon,\gamma_0,\gamma_1$, with $\Omega=\R^n$ and with the field
$(x,z,h)\mapsto A(h)$.  Its hypotheses are $n\ge2$, $p>1$, $\varepsilon\ge0$,
$\gamma_0>0$, $\gamma_1>0$ --- the last from $0<\gamma_0\le\gamma_1$ --- and
\noderef{EllipticStructureQ12} for that field, all hypotheses of the present
lemma.  Its conclusion holds for almost every $x$ and, the field being
independent of $(x,z)$, the inequality it asserts mentions neither; the
full-measure set of admissible $x$ is nonempty as in Step 1, so fixing one such
$x$ and any $z$ gives, for every $h\in\R^n$,
\[
|A(h)|\ \le\ |A(0)|+c_\ast\gamma_1\bigl(\varepsilon+|h|^2\bigr)^{\frac p2-1}|h| .
\tag{4}
\]
Since $|h|=(|h|^2)^{1/2}\le(\varepsilon+|h|^2)^{1/2}$ and the exponent $p/2-1$
acts on the same nonnegative base, (4) gives
\[
|A(h)|\ \le\ |A(0)|+c_\ast\gamma_1\bigl(\varepsilon+|h|^2\bigr)^{\frac{p-1}2} :
\tag{5}
\]
if $\varepsilon+|h|^2>0$ this is the identity
$(\varepsilon+|h|^2)^{\frac p2-1}(\varepsilon+|h|^2)^{1/2}
=(\varepsilon+|h|^2)^{\frac{p-1}2}$ for powers of a positive base, together
with the bound on $|h|$; and if $\varepsilon+|h|^2=0$ then $h=0$, so the second
summand on the right of (4) is $c_\ast\gamma_1\cdot0^{\frac p2-1}\cdot0=0$ and
the second summand on the right of (5) is
$c_\ast\gamma_1\cdot0^{\frac{p-1}2}=0$, the exponent $\frac{p-1}2$ being
positive because $p>1$ --- both right-hand sides equal $|A(0)|$, so (5) holds
in that case too.

Let $q'=p/(p-1)$, the H\"older conjugate of $p$, finite and $>1$ because
$p>1$.  We claim $|A(g_u)|\in L^{q'}(B_R)$.  Put
\[
\Phi=|A(0)|+c_\ast\gamma_1\bigl(\varepsilon+|g_u|^2\bigr)^{\frac{p-1}2}\ \ge 0 ,
\]
so that $|A(g_u)|\le\Phi$ pointwise on $B_R$, by (5) applied at
$h=g_u(x)$.  Both $|A(g_u)|$ and $\Phi$ are measurable by (3a), and
$t\mapsto t^{q'}$ is nondecreasing on $[0,\infty)$, so
$|A(g_u)|^{q'}\le\Phi^{q'}$ pointwise on $B_R$; by part (a) Domination of the
Preamble's proposition on the integral --- which is the part stated for a general
measure space, the monotonicity part (b) being stated only for interval
integrals --- it therefore suffices to show $\Phi\in L^{q'}(B_R)$, since then
$|A(g_u)|^{q'}$ is dominated by an integrable function and is itself
integrable.  By the $L^{q'}$ triangle
inequality (Minkowski, Preamble) it is enough to see that each of the two
summands of $\Phi$ is in $L^{q'}(B_R)$.  The first is
the constant $|A(0)|$, which is in $L^{q'}(B_R)$ because $|B_R|<\infty$.  For
the second, $\bigl((\varepsilon+|g_u|^2)^{\frac{p-1}2}\bigr)^{q'}
=(\varepsilon+|g_u|^2)^{\frac p2}$ since $\frac{p-1}2\cdot\frac p{p-1}
=\frac p2$, and pointwise
$\varepsilon+|g_u|^2\le2\max\{\varepsilon,|g_u|^2\}$, so raising to the power
$\frac p2>0$, under which $t\mapsto t^{p/2}$ is nondecreasing on $[0,\infty)$,
\[
\bigl(\varepsilon+|g_u|^2\bigr)^{\frac p2}
\le 2^{\frac p2}\max\bigl\{\varepsilon^{\frac p2},\,|g_u|^{p}\bigr\}
\le 2^{\frac p2}\bigl(\varepsilon^{\frac p2}+|g_u|^{p}\bigr),
\]
whose integral over $B_R$ is finite because
$\int_{B_R}\varepsilon^{p/2}=\varepsilon^{p/2}|B_R|<\infty$ and
$|g_u|\in L^p(B_R)$ by the $L^p$ clause of \noderef{MemW1pOn}.  Hence
$(\varepsilon+|g_u|^2)^{\frac{p-1}2}\in L^{q'}(B_R)$, so
$\Phi\in L^{q'}(B_R)$ and therefore
$|A(g_u)|\in L^{q'}(B_R)$.  The same argument, with $g_v$ in place of $g_u$,
gives $|A(g_v)|\in L^{q'}(B_R)$.

Consequently, by Cauchy--Schwarz pointwise (Preamble) and then H\"older with the
conjugate exponents $q'$ and $p$ on $B_R$ (Preamble) --- all three integrands
below being nonnegative and measurable, so that their integrals are well-defined
elements of $[0,\infty]$ and the chain is an inequality there ---
\[
\int_{B_R}\bigl|\bigl(A(g_u),g_\eta\bigr)\bigr|dx
\le\int_{B_R}|A(g_u)|\,|g_\eta|\,dx
\le\bigl\||A(g_u)|\bigr\|_{L^{q'}(B_R)}\bigl\||g_\eta|\bigr\|_{L^{p}(B_R)}<\infty ,
\]
using $|g_\eta|\in L^p(B_R)$ from (3); and likewise with $g_v$ in place of
$g_u$.  So both of the integrals appearing next are finite.

\emph{(3c) Testing and subtracting.}  Both $u$ and $v$ are weak solutions of
$\operatorname{div}A(\nabla\cdot)=0$ in $B_R$ (\noderef{IsWeakSolutionOn}, with
the operator $A(x,z,h)=A(h)$ and $B\equiv0$), so testing each against
$\varphi_\eta$, which is admissible by Step 2, gives
\[
\int_{B_R}\bigl(A(g_u),g_\eta\bigr)dx=0
\qquad\text{and}\qquad
\int_{B_R}\bigl(A(g_v),g_\eta\bigr)dx=0 .
\]
Subtracting --- legitimate because both integrals are finite by (3b) --- and
inserting (3),
\[
0=\int_{B_R}\bigl(A(g_u)-A(g_v),\,g_\eta\bigr)dx
=\int_{B_R}\bigl(A(g_u)-A(g_v),\,g_u-g_v\bigr)\chi_{\{w>\eta\}}\,dx .
\tag{6}
\]
Call the integrand of the right-hand side $F$.  It is measurable by (3a).  It is
integrable on $B_R$: by (3b) and the triangle inequality,
$|F|\le|A(g_u)-A(g_v)|\,|g_\eta|\le\bigl(|A(g_u)|+|A(g_v)|\bigr)|g_\eta|$, which
is integrable on $B_R$ by the two H\"older bounds of (3b).  And it is
nonnegative: on $B_R\cap\{w>\eta\}$ it equals
$(A(g_u)-A(g_v),g_u-g_v)$, which is $\ge0$ by (1) applied with $a=g_u(x)$ and
$b=g_v(x)$, while off that set it is $0$.  A nonnegative integrable function
with vanishing integral vanishes almost everywhere (Preamble), so $F=0$ a.e.\ on
$B_R$; in particular
\[
\bigl(A(g_u(x))-A(g_v(x)),\,g_u(x)-g_v(x)\bigr)=0
\qquad\text{for a.e.\ }x\in B_R\cap\{w>\eta\} ,
\]
and therefore, by (2), $g_u=g_v$ almost everywhere on $B_R\cap\{w>\eta\}$.  By
(3) this says
\[
g_\eta=0\qquad\text{almost everywhere on }B_R .
\tag{7}
\]

\emph{Step 4: a compactly supported function with vanishing gradient.}  By (3)
and (7), $\varphi_\eta\in W^{1,p}(B_R)$ has weak gradient $g_\eta$, and
$g_\eta=0$ almost everywhere on $B_R$.  We show first that $\varphi_\eta$ is
almost everywhere equal to a constant on $B_R$; this is where the hypothesis
\noderef{CitedSobolevPoincareBall} of the present lemma is consumed.

That hypothesis is stated for $W^{1,1}$ of the ball, so we record the passage
from $W^{1,p}$ to $W^{1,1}$.  The defining identity of
\noderef{HasWeakGradientOn} mentions no exponent, so it holds for the pair
$(\varphi_\eta,g_\eta)$ on $B_R$ exactly as given by (3).  For the memberships:
$\varphi_\eta\in L^p(B_R)$ and $|g_\eta|\in L^p(B_R)$ by (3), while
$|B_R|<\infty$ and $p\ge1$, so H\"older on a set of finite measure (Preamble)
puts both in $L^1(B_R)$.  Hence $\varphi_\eta\in W^{1,1}(B_R)$ with weak
gradient $g_\eta$ (\noderef{MemW1pOn}).

Let $C_6=C_6(n)>0$ be the constant of \noderef{CitedSobolevPoincareBall} and put
\[
c=\frac1{|B_R|}\int_{B_R}\varphi_\eta\,dx ,
\]
a finite real number because $\varphi_\eta\in L^1(B_R)$ and $0<|B_R|<\infty$.
Applying \noderef{CitedSobolevPoincareBall} to the ball $B_R(x_0)$, legitimate
because $R>0$ and $\varphi_\eta\in W^{1,1}(B_R)$ with weak gradient $g_\eta$,
\[
\Bigl(\int_{B_R}|\varphi_\eta-c|^{\frac n{n-1}}dx\Bigr)^{\frac{n-1}n}
\ \le\ C_6\int_{B_R}|g_\eta|\,dx\ =\ 0 ,
\]
the right-hand side vanishing because $g_\eta=0$ almost everywhere on $B_R$ by
(7), so that its integral is $0$.  Write
$X=\int_{B_R}|\varphi_\eta-c|^{n/(n-1)}dx$, the integral of a nonnegative
measurable function and hence a well-defined element of $[0,\infty]$.  The
exponent $\frac{n-1}n$ is positive because $n\ge2$, so $X^{\frac{n-1}n}\ge0$;
with the display this forces $X^{\frac{n-1}n}=0$, and since $X>0$ would give
$X^{\frac{n-1}n}>0$ we conclude $X=0$.  In particular
$|\varphi_\eta-c|^{n/(n-1)}$ is integrable on $B_R$ with integral $0$, so, being
nonnegative, it vanishes almost everywhere on $B_R$ (Preamble); and since
$\frac n{n-1}>0$, a vanishing $|\varphi_\eta-c|^{n/(n-1)}$ forces
$|\varphi_\eta-c|=0$.  Therefore
\[
\varphi_\eta=c\qquad\text{almost everywhere in }B_R .
\]

By Step 2 the support of $\varphi_\eta$ is contained in
the compact set $K_\eta\subseteq B_R$, so $B_R\setminus K_\eta$ is open, being
the intersection of the open set $B_R$ with the complement of the closed set
$K_\eta$, and it is nonempty: were it empty we would have $B_R=K_\eta$, making
the open ball $B_R$ compact, which it is not.  A nonempty open subset of $\R^n$
has positive measure, and $\varphi_\eta=0$ on it.  Hence $c=0$, i.e.
\[
\varphi_\eta=0\quad\text{a.e.\ in }B_R,
\qquad\text{that is}\qquad
u-v\le\eta\quad\text{a.e.\ in }B_R .
\tag{8}
\]

\emph{Step 5: letting $\eta\downarrow0$.}  Apply (8) with $\eta=1/j$ for
$j=1,2,3,\dots$: for each $j$ there is a null set $N_j\subseteq B_R$ outside of
which $u-v\le1/j$.  The union $N=\bigcup_jN_j$ is null, and for every
$x\in B_R\setminus N$ we have $u(x)-v(x)\le1/j$ for every $j$, hence
$u(x)-v(x)\le0$.  Therefore $u\le v$ almost everywhere in $B_R$, which is the
assertion.
\end{proof}
